\documentclass[12pt]{article}

\usepackage{amsmath, amssymb, amsthm}
\usepackage{enumerate}
\usepackage{hyperref}
\usepackage{geometry}
\usepackage{newtxtext}
\usepackage{newtxmath}
\geometry{margin=1in}

\title{Introduction to Functional Analysis - Lecture 1}
\author{}
\date{}

% ----------------------------------------------------------------
\begin{document}
\maketitle

\textit{Young's inequality} appears throughout real analysis, the theory of $L^p$ spaces,
and partial differential equations, and it is a key step in proving the celebrated H\"{o}lder inequality.

In this problem, we will prove it algebraically using single-variable calculus following a proof in Casella and Berger's book on Statistical Inference (wow!). This is one of the more elegant proofs I know, and it highlights how calculus is useful. The exposition follows that from a video by Robert Cruikshank.

\medskip\noindent
\textbf{Setup.}
Throughout, let $p$ and $q$ be real numbers greater than $1$ satisfying
\[
  \frac{1}{p}+\frac{1}{q}=1.
\]
Such a pair $(p,q)$ is called a pair of \textit{conjugate exponents}. For example, consider $(p,q)=(2,2)$.

\section*{Questions}

\begin{enumerate}[(a)]


% -------- algebraic lemma ------------------------------------
\item Using the condition $\frac{1}{p}+\frac{1}{q}=1$, prove by algebraic manipulations that
\[
  (p-1)q = p.
\]
Deduce that $a^{(p-1)q} = a^p$ using $a > 0$ as hypothesis.

% -------- algebraic proof: setup ----------------------------
\item Fix $b>0$ (without loss of generality, why?) and define $f\colon[0,\infty)\to\mathbb{R}$ by
\[
  f(a) = \frac{a^p}{p}+\frac{b^q}{q}-ab.
\]
Compute $f'(a)$ for $a>0$, and describe the critical point $a_0\in(0,\infty)$.

% -------- global minimum ------------------------------------
\item Show that $f$ attains its global minimum on $[0,\infty)$ at $a_0$. (Hint: use the second derivative test, do not forget to extend this to the global case.)

% -------- minimum value is zero -----------------------------
\item Use part (a) to compute $f(a_0)$, and show that $f(a_0)=0$.
      Deduce \textit{Young's inequality}: for all $a,b\geq 0$,
      \[
        ab \leq \frac{a^p}{p}+\frac{b^q}{q}.
      \]
\end{enumerate}

% ================================================================
\newpage
\section*{Solutions}

\begin{enumerate}[(a)]

\item We compute
      \[
        f'(a) = a^{p-1}-b.
      \]
      Setting $f'(a_0)=0$ gives $a_0^{p-1}=b$, so
      \[
        a_0 = b^{1/(p-1)}.
      \]
      This is the unique positive solution because $a\mapsto a^{p-1}$ is
      strictly increasing on $(0,\infty)$.

\item For $a < a_0$ we have $a^{p-1} < b$, so $f'(a)<0$ and $f$ is decreasing.
      For $a > a_0$ we have $a^{p-1}>b$, so $f'(a)>0$ and $f$ is increasing.
      Hence $f$ decreases up to $a_0$ and then increases, so $f(a_0)$ is a global minimum on $[0,\infty)$.

      (Alternatively: $f''(a)=(p-1)a^{p-2}>0$ for all $a>0$,
      so $f$ is strictly convex, and any interior critical point is a global minimum.)

\item Using $a_0 = b^{1/(p-1)}$ and the identity $(p-1)q=p$ from part (a),
      we get $\dfrac{p}{p-1}=q$, so
      \[
        a_0^p = \bigl(b^{1/(p-1)}\bigr)^p = b^{p/(p-1)} = b^q.
      \]
      Likewise,
      \[
        a_0\,b = b^{1/(p-1)}\cdot b = b^{1/(p-1)+1} = b^{p/(p-1)} = b^q.
      \]
      Therefore
      \[
        f(a_0)
        = \frac{a_0^p}{p}+\frac{b^q}{q}-a_0\,b
        = \frac{b^q}{p}+\frac{b^q}{q}-b^q
        = b^q\!\left(\frac{1}{p}+\frac{1}{q}-1\right)
        = b^q\cdot 0
        = 0.
      \]
      Since $a_0$ is the global minimum of $f$, we have $f(a)\geq f(a_0)=0$ for all $a\geq 0$,
      that is,
      \[
        \frac{a^p}{p}+\frac{b^q}{q}-ab\geq 0.
      \]
      Rearranging gives Young's inequality $ab\leq\dfrac{a^p}{p}+\dfrac{b^q}{q}$.

\end{enumerate}

\end{document}